\documentclass[preprint]{glimmer-paper}

\title{Reproducibility as Knowledge Graph Navigation:\\
A Research-Object Knowledge Base for AI-Native Neuroimaging Analysis}
\author{Shady El Damaty}
\affil{Holonym / Opscientia}
\version{0.5.1}
\repo{https://github.com/hebbianloop/glimmer}
\date{2026}

\newcommand{\code}[1]{\texttt{#1}}
\newcommand{\node}[1]{\textsf{#1}}

\begin{document}
\maketitle

\begin{abstract}
Reproducible-pipeline tooling---BIDS for data layout, DataLad and git-annex for
content-addressed versioning, Nipype and containers for pinned computation---made the
\emph{bytes} of a neuroimaging analysis recoverable and gave projects syntactic
structure. What it did not give is a \emph{navigable knowledge layer}: a typed,
machine-traversable account of which datasets, methods, derivatives, findings, and
publications a project comprises, how they relate, and what evidence grounds each claim.
This gap matters acutely for AI-native workflows, where an agent must reason over a
project's provenance to render decisions a human can audit. We present Glimmer, a
research-object knowledge base that models a research project as a typed-entity graph
distributed across per-entity sidecar files. Twelve core node types and a versioned edge
taxonomy turn datasets, methods, derivatives, findings, standards, publications, research
questions, people, and programs into first-class nodes; every node carries a content hash
and its DataLad re-fetch coordinates, and every agent-produced node carries a mandatory
reasoning trace citing the nodes it read. A schema-level verifiability contract
distinguishes \emph{exact} provenance (a deterministic result must re-run from its cited
SHAs) from \emph{structural} provenance (an inferred claim must cite real nodes that
contain the values it reports), and a domain-profile mechanism keeps the core
domain-neutral while letting neuroimaging (BIDS) and other fields add their own
vocabulary. Because the graph is plain files, it survives \code{git clone},
\code{datalad export}, and \code{rsync} with no database. We describe the schema, the
agent protocol, the agentic research loop it enables, and two worked examples (a
DataLad$\rightarrow$Nipype pipeline and a literature-retrieval adapter), and we frame
reproducibility as a property an agent establishes by \emph{navigating} the graph and
re-running from cited coordinates. Turning that navigable property into an automatic
execution-and-verification engine is the subject of follow-on work.
\end{abstract}

\section{Introduction}

A decade of investment in reproducible pipelines transformed neuroimaging. The Brain
Imaging Data Structure \citep{gorgolewski2016bids} standardized how raw data is laid out;
DataLad and git-annex \citep{halchenko2021datalad} made datasets content-addressed and
re-fetchable by hash; Nipype \citep{gorgolewski2011nipype} and containerized BIDS-Apps
such as MRIQC \citep{esteban2017mriqc} and fMRIPrep \citep{esteban2019fmriprep} pinned the
toolchain so a computation could, in principle, be re-executed. Together these guarantee
two things: the \emph{bytes} an analysis consumed can be recovered, and the project has a
\emph{syntactic} structure a tool can parse.

They do not guarantee a third thing that matters increasingly: a \emph{navigable knowledge
layer}. Given a published number, what dataset produced it, under which method version,
conforming to which standard, supporting which claim, cited in which paper, framed against
which hypothesis? In practice that knowledge lives in lab wikis, file-naming conventions,
analysis scripts, and researchers' memories---none of it typed, none of it traversable,
and none of it survives the people who built it. The reproducibility crisis is, in this
framing, partly a \emph{knowledge-representation} problem: the relationships among a
project's artifacts are real but unrecorded.

The problem sharpens under AI-native workflows. An agent asked to quality-control a
dataset, summarize a result, or draft a paragraph must reason over exactly these
relationships, and---if its output is to be trusted---must leave behind an auditable
account of the evidence it used. A free-text answer from a language model is not
auditable; a decision grounded in a typed graph, with an explicit trace of which nodes
were read and which metrics were treated as load-bearing, is.

\paragraph{Thesis.} We propose to model a research project as a typed-entity
\emph{research-object graph}, distributed across per-entity sidecar files that live
alongside the data in version control, so that reproducibility becomes
\emph{knowledge-graph navigation}: an agent (or a human) answers a provenance question by
traversing typed edges from a claim back to the data and code that ground it, and
establishes that a result reproduces by re-fetching and re-running from the coordinates
the graph records. Glimmer is the schema, the verifiability contract, and a reference
implementation of this idea. This paper describes the architecture as of v0.5.x; an
automatic execution-and-verification engine that turns the navigable property into a
machine-checked one is forthcoming \citep{eldamaty2026proactive}.

\paragraph{Contributions.}
\begin{itemize}
  \item A domain-neutral schema of twelve core node types and a versioned edge taxonomy
  that represents the data layer, the evidence layer, and the social/research-program
  layer of a project as one traversable graph (\S\ref{sec:schema}).
  \item A distributed, file-based realization---per-entity YAML sidecars plus a root
  index---that carries content hashes and DataLad re-fetch coordinates on every node and
  survives ordinary distribution tools (\S\ref{sec:schema}, \S\ref{sec:discussion}).
  \item An agent protocol: a schema-level verifiability contract requiring every
  agent-produced node to declare its identity and emit a reasoning trace, and
  distinguishing exact from structural verification (\S\ref{sec:protocol}).
  \item The agentic research loop the substrate enables---plans-as-issues, concept
  decomposition, and four specialized agent roles---and a retrieval adapter that gives
  literature search the same provenance discipline (\S\ref{sec:loop}).
  \item Two worked examples and a domain-profile mechanism that keeps the core
  domain-neutral while supporting BIDS-specific vocabulary (\S\ref{sec:examples}).
\end{itemize}

\section{Background and related work}
\label{sec:background}

\paragraph{Data and pipeline standards.} BIDS \citep{gorgolewski2016bids} fixes a
directory layout and sidecar conventions for neuroimaging; it is syntactic structure, not
a relational account of a project. DataLad and git-annex \citep{halchenko2021datalad}
provide content-addressed, distributed version control: a dataset is a git repository
whose large files are tracked by annex key, so any file can be re-fetched by a pinned
commit and key. Nipype \citep{gorgolewski2011nipype} wraps heterogeneous tools (FSL,
FreeSurfer, ANTs) behind a uniform interface and records workflow structure. Glimmer does
not replace any of these; it reads them and sits above them, treating a DataLad
superdataset as the source of truth and recording, on each node, the coordinates needed to
recover the bytes.

\paragraph{Provenance description.} The W3C PROV data model \citep{w3cprov2013}
formalizes entities, activities, and agents; the Neuroimaging Data Model (NIDM)
\citep{maumet2016nidm} applies a PROV/RDF account to statistical results; RO-Crate
\citep{soilandreyes2022rocrate} packages a dataset with typed, linked-data metadata; and
schema.org provides a general vocabulary for web-discoverable description. These are
description formats. Glimmer is deliberately compatible with them---its node types map
cleanly onto PROV entities and activities, and a Glimmer graph can be cross-read as
RO-Crate or JSON-LD---but its contribution is a \emph{working substrate} an agent
navigates and a validator enforces, not a new serialization. The interoperability stance
is explicit: format does not matter if an agent can translate between formats, so a
sidecar may be BIDS-native JSON, Glimmer-native YAML, or any structured form the agent can
parse into the same typed graph.

\paragraph{AI for science.} End-to-end autoresearch systems
\citep{lu2024aiscientist,boiko2023coscientist} run fixed internal pipelines from idea to
manuscript. Glimmer is complementary: rather than a pipeline, it is a typed, versioned,
auditable substrate that multiple agents---possibly from different systems---can reason
over and contribute to, with a verifiability contract that makes each contribution
checkable. The graph survives independently of any one autoresearch system.

\section{The Glimmer schema}
\label{sec:schema}

A Glimmer research-object knowledge base (RO-KB) is a set of typed nodes, each stored as a
sidecar file, enumerated by a root index. We describe the node types, the edges, the
sidecar format, and the domain-profile mechanism.

\subsection{Node types}

The core schema is domain-neutral: vocabulary fixed by a particular field's standards
lives in profiles (\S\ref{sec:profiles}), not in the core. Twelve node types span three
layers.

\emph{The data and evidence layer.}
\begin{itemize}
  \item \node{dataset} --- research data of any kind. The core carries only identity,
  generic provenance, and re-fetch coordinates; kind-specific attributes
  (participant, session, modality) come from a profile.
  \item \node{method} --- a named analysis tool, pipeline, or workflow (an FSL binary, a
  Nipype interface or workflow, a script), with \code{tool}, \code{version},
  \code{parameters}, a \code{parameters-hash}, and a \code{workflow-definition-sha}.
  \item \node{experiment} --- a task or acquisition \emph{paradigm} (conditions, stimuli,
  timing), distinct from a static standard.
  \item \node{instrument} --- the measurement apparatus that \emph{generates} data (a
  scanner, survey, assay, task-delivery system, or device); a dataset is
  \code{acquired-with} an instrument.
  \item \node{derivative} --- the output of applying a method to a dataset; a first-class
  node, not a directory, carrying an \code{output-hash} so a re-run can be checked, and a
  required \code{provenance-mode} (\S\ref{sec:protocol}).
  \item \node{finding} --- an interpreted assertion grounded in one or more derivatives,
  the unit between ``the pipeline produced this output'' and ``we wrote a paper about
  it,'' aligned with evidence-graph practice: an interpretation plus evidence pointers
  plus verifiable provenance.
  \item \node{standard} --- a specification, atlas, template, or protocol (a BIDS version,
  an atlas, a QC rating scale), modeled as a node so constraints are edges an agent can
  read.
  \item \node{publication} --- a paper draft, abstract, or preprint that aggregates
  findings.
\end{itemize}

\emph{The social and research-program layer.}
\begin{itemize}
  \item \node{concept} --- a research question, hypothesis, or theme: the unit a research
  program operates at (what a grant funds, what a thesis defends). The agentic loop
  decomposes a concept into sub-hypotheses.
  \item \node{persona} --- a person or organizational role; the in-graph identity an
  attribution edge resolves to.
  \item \node{organization} --- an institution, lab, consortium, journal, or funder.
  \item \node{program} --- a study, cohort, or initiative as a first-class container that
  organizes datasets, experiments, concepts, and publications around one mission; programs
  nest via \code{part-of} to form multi-study hierarchies.
\end{itemize}

\subsection{Edges}

Edges are properties on the source node, each carrying a type and a target. The taxonomy
has three families.

\emph{Structural edges} record the data-to-claim lineage: \code{produced-by}
(derivative/dataset $\rightarrow$ method), \code{derives-from} (derivative $\rightarrow$
dataset or upstream derivative), \code{applies-to} and \code{produces} (method
$\leftrightarrow$ derivative), \code{conforms-to} (\textrm{$\rightarrow$} standard), the
\code{cites-*} family (publication $\rightarrow$ dataset/method/derivative/finding), the
method-pipeline DAG (\code{composes}, \code{upstream-of}, \code{downstream-of}), and the
two acquisition edges promoted to the core in v0.5, \code{acquired-with} (dataset or
experiment $\rightarrow$ instrument) and \code{described-by} (dataset $\rightarrow$
experiment).

\emph{Meta-graph edges} record evidence and the research program: \code{based-on}
(finding $\rightarrow$ derivative/dataset), \code{addresses-concept} (finding or
publication $\rightarrow$ concept), \code{tests-hypothesis} (experiment $\rightarrow$
concept), the concept-relation edges \code{decomposes-into}, \code{extends-concept},
\code{subsumed-by}, \code{competes-with}, and \code{superseded-by}, the evidence-relation
edges \code{supports} and \code{contradicts}/\code{challenged-by}, and the in-graph
attribution layer \code{authored-by}, \code{affiliated-with}, \code{funded-by},
\code{mentors}, \code{leads}, and \code{part-of}.

\emph{Universal edges} are allowed from any node and unioned in by the validator:
\code{contributed-by} (an out-of-graph contributor identifier---an ORCID URI, email, or
kebab id---carrying role metadata), \code{in-program} (membership in a program in this
graph), and \code{cross-project} (a relationship to a node in \emph{another} project's
graph, addressed by a namespaced id such as \code{ads-glimmer:org-nij}). The last two let
a subproject inherit a canonical node from its parent program, or a harmonizing analysis
connect two graphs, without a federated index.

\subsection{Sidecars, the index, and provenance}

Each node is a Markdown file with YAML front-matter; standalone sidecars are
Glimmer-native YAML, while a sidecar extending a BIDS file in place is BIDS-native JSON
augmented with an \code{\_x-glimmer} block. Every node has \code{id}, \code{type},
\code{name}, \code{created}, \code{modified}, an optional \code{description}, an optional
\code{edges} list, and a \code{provenance-hash}: a SHA-256 over the node's body, so a git
commit is a snapshot of the entire graph's state. A root file, \code{\_glimmer-index.json},
enumerates every node id and its path; it is the mandatory load for an agent.

Re-fetchability is carried on the node. Datasets, derivatives, and methods record DataLad
coordinates---\code{datalad-superdataset}, \code{datalad-relative-path},
\code{datalad-commit-sha}, and \code{datalad-annex-key}---so the graph is self-describing
for recovery: installing the superdataset and getting the cited path reproduces exactly
the bytes the analysis ran on. Version 0.5.1 adds a \emph{data-availability} block to
\node{dataset} (\code{data-remote}, \code{data-provenance}, \code{data-last-commit},
\code{data-liveness}) so a query over the graph can tell whether a dataset is currently
pullable, and from which storage backend, complementing the re-fetch coordinates.

\subsection{Domain profiles}
\label{sec:profiles}

Keeping the core domain-neutral requires a disciplined extension mechanism. A
\emph{domain profile} is a small YAML file that augments core node types with extra
required or optional fields and, as of v0.5, may declare domain-specific node types and
edge types. Profiles come in two tiers: a curated library versioned in the repository
(\code{glimmer/schema/profiles/<domain>.yaml}) and a researcher's own, local to one
knowledge base (\code{<rokb>/\_glimmer-profiles/<domain>.yaml}); a node selects its profile
by its \code{domain} field, falling back to a KB-level default. The shipped neuroimaging
profile encodes BIDS vocabulary---augmenting \node{dataset} with \code{subject-id},
\code{session}, \code{modality}, \code{scanner}, and \code{bids-version}; \node{experiment}
with \code{design} and \code{stimulus-set}; \node{method} with a Nipype node kind; and
\node{derivative} with an \code{output-kind}---and declares the study-provenance edges a
neuroimaging project needs. The governance rule is simple: a node or edge type is promoted
to the core only once at least two domains need it; otherwise it stays in a profile. This
lets a new field adopt Glimmer without forking the schema.

\section{The agent protocol: a verifiability contract}
\label{sec:protocol}

The schema records structure; the agent protocol makes agent contributions auditable. Its
rules are enforced by the validator, so a graph that violates them does not validate.

\paragraph{Identity and reasoning trace.} When a \node{finding} or \node{derivative} is
produced by an agent rather than by a deterministic computation, it must set
\code{produced-by-agent} (a stable model identifier) and carry a \code{reasoning-trace}.
The trace is structured: \code{nodes-accessed} lists every node id the agent read,
\code{metrics-cited} records the numeric or categorical evidence treated as load-bearing,
\code{evidence-summary} is a short justification with inline node-id citations, and
\code{model-identifier} and \code{timestamp} fix provenance. A finding must additionally
carry a \code{based-on} chain naming the derivatives or datasets that support it; a finding
without one is treated as ungrounded.

\paragraph{Modes of provenance.} Every \node{derivative} declares a \code{provenance-mode}:
\code{deterministic} (a Nipype/FSL computation whose output is fixed by its inputs),
\code{agent-inferred} (a language-model summary), or \code{stochastic} (e.g.\ a randomized
initialization or a retrieval over an embedding index). The mode tells a reader what kind
of verification is even possible.

\paragraph{Exact versus structural verification.} The contract distinguishes two regimes.
For a \code{deterministic} output, verification is \emph{exact} in principle: re-running
the cited method at the cited input SHAs must reproduce the recorded \code{output-hash}, so
the graph's claim about the artifact is falsifiable by re-execution. For an
\code{agent-inferred} output, verification is \emph{structural}: the reasoning trace must
cite nodes that exist, those nodes must contain the values the trace reports, and the
interpretation must be a plausible reading of them. Neither test guarantees correctness;
together they guarantee \emph{auditability}---a reviewer can walk from a claim back through
its cited evidence and check that the evidence exists, conforms to the standards it cites,
and says what the agent says it says. In this version these regimes are a contract the
validator enforces and a discipline a human or a CI job applies by re-running; an engine
that performs the re-execution and the structural check automatically is the subject of
follow-on work (\S\ref{sec:limits}).

\paragraph{Sanity-checking against source.} A method records a
\code{workflow-definition-sha} and may record a \code{source-checkout-url} and
\code{source-build-instructions}. When a reasoning trace cites a derivative whose method
version has drifted from upstream source, the recovery move---borrowed from long-running
neuroimaging projects---is to rebuild the method from source rather than trust a possibly
divergent cached binary. The schema makes this expressible; an agent may choose to enforce
it.

\section{The agentic research loop}
\label{sec:loop}

The substrate is designed for a specific operating pattern: \emph{plans-as-issues} applied
to research. An author writes a research idea as a \node{concept} node with a tracked
issue; the concept is decomposed via \code{decomposes-into} into falsifiable
sub-hypotheses; and for each sub-hypothesis the project runs one or more agent loops whose
outputs route to human review (approve, revise, or reject). Four agent roles recur, each
operating with a different access scope and the same verifiability contract:

\begin{itemize}
  \item a \emph{literature scout} that traverses external bibliographic sources and the
  project's existing publications and emits new \node{publication} nodes for relevant prior
  work, read-only on local data;
  \item a \emph{QC agent} that rates datasets and their derivatives against a standard,
  emitting artifacts with full reasoning traces;
  \item an \emph{analysis agent} that runs deterministic computation and emits
  \node{derivative} nodes with \code{provenance-mode: deterministic}; and
  \item a \emph{synthesis agent} that walks the populated graph and drafts a
  \node{publication} by composing derivatives, QC artifacts, and cited publications into an
  argument, with its output gated on human review.
\end{itemize}

At every step the output's reasoning trace cites the standards it applied, so a reviewer
can verify standards adherence without inspecting the data directly; the protocol makes
verifiability non-optional at each step. Crucially, the loop is the \emph{operational
pattern the substrate enables}---typed roles emitting typed, traceable nodes---rather than
a monolithic executable; the substrate's job is to make every contribution navigable and
checkable.

\paragraph{Retrieval with provenance.} Literature search is inherently non-deterministic,
so the retrieval adapter gives it the honest analogue of SHA-based reproducibility. A
retrieval-grounded \node{finding} is marked \code{provenance-mode: stochastic} and embeds a
\code{retriever-manifest}---the embedding model and version, the chunking and reranking
configuration, and the index SHA---so a same-configuration re-run is reproducible modulo
the index. The adapter's indexing operation is idempotent on a node's id and
provenance-hash, so re-emitting an unchanged node is a no-op. This keeps even the
stochastic, external-knowledge parts of the loop inside the same auditable contract.

\section{Worked examples}
\label{sec:examples}

\paragraph{A DataLad$\rightarrow$Nipype pipeline.} The reference example operates on the
public OpenNeuro dataset \code{ds000114}. A short install step runs
\code{datalad install ///openneuro/ds000114} and selectively gets a subject's T1-weighted
image, recording the DataLad coordinates. A Nipype workflow applies skull-stripping (FSL
BET) and tissue segmentation (FSL FAST); the emitter then writes typed sidecars: a
\node{dataset} for the input, \node{method} nodes for BET and FAST (each with its tool
version, parameter hash, and workflow-definition SHA), \node{derivative} nodes for the
outputs (each with an \code{output-hash} and \code{provenance-mode: deterministic}), and a
\node{finding} reporting a derived brain volume. The graph is self-describing for recovery:
the cited superdataset, commit, and annex key locate exactly the input bytes, and the
method's SHA plus parameters specify the computation, so the derivative's recorded hash is
the target a re-run must reproduce. The validator confirms that every edge target exists,
that the finding carries a \code{based-on} chain, and that the reasoning-trace contract
holds.

\paragraph{A literature-retrieval adapter.} The second example models the literature-scout
role. A \node{concept} states a falsifiable hypothesis (a naturalistic emotional-film fMRI
paradigm predicting a later behavioral outcome); the scout retrieves prior work, emits
\node{publication} nodes with citation metadata, and emits a \node{finding} whose
\code{reasoning-trace} records the retrieval scores and the publications it read, marked
\code{provenance-mode: stochastic} with its \code{retriever-manifest}. A reviewer can read
the trace, follow \code{based-on} to each cited publication, and confirm the finding cites
real, retrievable evidence---structural verification in practice.

\section{Discussion}
\label{sec:discussion}

\paragraph{Distributed-first by design.} Glimmer deliberately avoids a central database.
Because the graph is per-entity files committed alongside the data, it survives the tools
researchers already use---\code{git clone}, \code{datalad export}, \code{rsync}---with no
bespoke serialization, and a single commit at the superdataset level captures the full
state of data, code, derivatives, and graph together. The cost---no global query without
loading the index---is acceptable for project- and program-scale graphs and buys
durability and tool-independence.

\paragraph{Domain-agnostic core, domain-specific profiles.} The split between a
twelve-type core and curated/local profiles lets the same architecture serve neuroimaging
today and other compute-intensive, standards-backed fields tomorrow, with promotion to the
core gated on cross-domain demand. The reference implementation is neuroimaging because
that is where the standards ecosystem is most mature, but nothing in the core is
neuroimaging-specific.

\paragraph{A junction with decentralized science.} Because nodes carry stable, possibly
out-of-graph identities (ORCID, namespaced cross-project ids) and the substrate is
file-based and content-addressed, Glimmer composes naturally with decentralized-science
infrastructure---OpenNeuro and Opscientia for data, and verifiable-researcher-identity
systems for attribution and signing. The \code{cross-project} edge is the seam along which
independent project graphs federate.

\section{Limitations and future work}
\label{sec:limits}

The central limitation of this version is that verification is a \emph{navigable property},
not an automatic one. The schema records everything needed to reproduce a deterministic
result---input SHAs, method SHA, parameters, the target output hash---and the validator
enforces the reasoning-trace contract, but the act of re-fetching, re-running, and
comparing is performed by a human or a continuous-integration job, not by the substrate
itself. An \emph{executable node runner} that gates a computation's inputs, replays it in
its pinned environment, and verifies its outputs automatically---turning the exact and
structural regimes of \S\ref{sec:protocol} into machine-checked verdicts---is the natural
next step and the subject of follow-on work \citep{eldamaty2026proactive}. Other open
directions include federation across institution graphs with a reputation model for schema
extensions, durable-storage guarantees expressed in the model rather than left to
operations, and temporal edges for longitudinal study waves.

\section{Conclusion}

Reproducible pipelines gave neuroimaging recoverable bytes and syntactic structure;
Glimmer adds the layer above them, a typed research-object graph in which datasets,
methods, derivatives, findings, standards, publications, questions, and people are
first-class nodes connected by versioned edges, distributed across files that survive
ordinary tooling. By carrying content hashes and DataLad coordinates on every node and a
mandatory reasoning trace on every agent output, the graph makes reproducibility a matter
of \emph{navigation}: an agent or reviewer walks from a claim to the evidence that grounds
it and re-runs from the coordinates the graph records. That navigable foundation is what an
automatic execution-and-verification engine can then build upon.

\bibliographystyle{plainnat}
\bibliography{refs}

\end{document}
